\section{Численные результаты}
\label{sec:Chapter5} \index{Chapter5}

Вот результаты для $\mathsf{LAC}$ и $\mathsf{NewHope}$. Параметры $[n, q, k]$ были взяты $[1024, 251, 1]$ и $[1024, 12289, 8]$ соответственно. Целевой DFR положим $10^{-6}$.

Для $\mathsf{LAC}$ имеет смысл рассматривать размер алфавита $Q$ от $2$ до $6$ символов включительно. Посчитав распределение в $\Z_q$ шума на компьютере, бинпоиском подбираем необходимое кол-во исправляемых ошибок.

\begin{table}[h]
	\centering
	\begin{minipage}{0.48\textwidth}
		\centering
		\begin{tabular}{|c|c|c|}
			\hline Q & t & err \\ \hline
			2 & 20 & 0.0054 \\ \hline
			3 & 105 & 0.0636 \\ \hline
			4 & 227 & 0.1640 \\ \hline
			5 & 334 & 0.2599 \\ \hline
			6 & 436 & 0.3537 \\ \hline
		\end{tabular}
		\caption*{Hamming}
	\end{minipage}
	\begin{minipage}{0.48\textwidth}
		\centering
		\begin{tabular}{|c|c|c|c|}
			\hline Q & t & err1 & err2 \\ \hline
			2 & 20 & 0.0054 & 0 \\ \hline
			3 & 105 & 0.0636 & 0 \\ \hline
			4 & 227 & 0.1640 & 3.32e-05 \\ \hline
			5 & 336 & 0.2591 & 0.0009 \\ \hline
			6 & 444 & 0.3478 & 0.0058 \\ \hline
		\end{tabular}
		\caption*{Li}
	\end{minipage}
\end{table}

Последний столбец отвечает за вероятность ошибки в каждом символе (если рассматривать их независимо). В метрике Ли это вероятности отличия от исходного символа на расстояние $1$ и $2$ соответственно.

Для $\mathsf{NewHope}$ делаем то же самое, только рассматривая другие $Q$.

\begin{table}[h]
	\centering
	\begin{tabular}{|c|c|c|}
		\hline Q & t & err \\ \hline
		2 & 0 & $\approx 0$ \\ \hline
		3 & 0 & $\approx 0$ \\ \hline
		4 & 0 & $\approx 0$ \\ \hline
		5 & 0 & 1.50e-11 \\ \hline
		6 & 1 & 1.68e-08 \\ \hline
		7 & 1 & 1.30e-06 \\ \hline
		8 & 3 & 2.25e-05 \\ \hline
		9 & 4 & 1.64e-04 \\ \hline
	\end{tabular}
	\caption*{Hamming, Li}
	
	\begin{minipage}{0.48\textwidth}
		\centering
		\begin{tabular}{|c|c|c|}
			\hline Q & t & err \\ \hline
			10 & 7 & 0.0007 \\ \hline
			12 & 18 & 0.0047 \\ \hline
			14 & 38 & 0.0153 \\ \hline
			16 & 65 & 0.0337 \\ \hline
			18 & 100 & 0.0593 \\ \hline
			20 & 138 & 0.0894 \\ \hline
			23 & 198 & 0.1395 \\ \hline
			26 & 258 & 0.1914 \\ \hline
			30 & 332 & 0.2574 \\ \hline
			34 & 397 & 0.3173 \\ \hline
			38 & 454 & 0.3707 \\ \hline
		\end{tabular}
		\caption*{Hamming}
	\end{minipage}
	\hfill
	\begin{minipage}{0.48\textwidth}
		\centering
		\begin{tabular}{|c|c|c|c|}
			\hline Q & t & err1 & err2 \\ \hline
			10 & 7 & 0.0007 & $\approx 0$ \\ \hline
			12 & 18 & 0.0047 & $\approx 0$ \\ \hline
			14 & 38 & 0.0153 & $\approx 0$ \\ \hline
			16 & 65 & 0.0337 & 2.28e-10 \\ \hline
			18 & 100 & 0.0593 & 1.71e-08 \\ \hline
			20 & 138 & 0.0894 & 3.81e-07 \\ \hline
			23 & 198 & 0.1395 & 9.89e-06 \\ \hline
			26 & 258 & 0.1913 & 0.0001 \\ \hline
			30 & 333 & 0.2567 & 0.0007 \\ \hline
			34 & 401 & 0.3145 & 0.0028 \\ \hline
			38 & 464 & 0.3633 & 0.0074 \\ \hline
		\end{tabular}
		\caption*{Li}
	\end{minipage}
\end{table}

Теперь мы знаем минимальные расстояния между кодами $d = 2t+1$.

Зная их и $n$, посчитаем границы Хэмминга и Варшамова-Гилберта на размер кода $M$ и оценим его скорость $R = \frac{\log_2(M)}n$ в $Q$-ичном алфавите.

\begin{table}[h]
	\centering
	\begin{minipage}{0.48\textwidth}
		\centering
		\begin{tabular}{|c|c|c|}
			\hline Q & $R_{GV}$ & $R_H$ \\ \hline
			2 & 0.7658 & 0.8646 \\ \hline
			3 & 0.6526 & 1.0098 \\ \hline
			4 & 0.3113 & 0.8902 \\ \hline
			5 & 0.0895 & 0.7635 \\ \hline
			6 & 0.0010 & 0.6172 \\ \hline
		\end{tabular}
		\caption*{LAC, Hamming}
	\end{minipage}
	\begin{minipage}{0.48\textwidth}
		\centering
		\begin{tabular}{|c|c|c|}
			\hline Q & $R_{GV}$ & $R_H$ \\ \hline
			2 & 0.7658 & 0.8646 \\ \hline
			3 & 0.6526 & 1.0098 \\ \hline
			4 & 0.4785 & 1.0000 \\ \hline
			5 & 0.3761 & 1.0050 \\ \hline
			6 & 0.3319 & 1.0131 \\ \hline
		\end{tabular}
		\caption*{LAC, Li}
	\end{minipage}
	
	\bigskip
	
	\begin{minipage}{0.48\textwidth}
		\centering
		\begin{tabular}{|c|c|c|}
			\hline Q & $R_{GV}$ & $R_H$ \\ \hline
			2 & 1.0000 & 1.0000 \\ \hline
			3 & 1.5850 & 1.5850 \\ \hline
			4 & 2.0000 & 2.0000 \\ \hline
			5 & 2.3219 & 2.3219 \\ \hline
			6 & 2.5619 & 2.5729 \\ \hline
			7 & 2.7838 & 2.7951 \\ \hline
			8 & 2.9342 & 2.9650 \\ \hline
			9 & 3.0833 & 3.1236 \\ \hline
			10 & 3.1775 & 3.2439 \\ \hline
			12 & 3.2475 & 3.3999 \\ \hline
			14 & 3.1555 & 3.4450 \\ \hline
			16 & 2.9596 & 3.4151 \\ \hline
			18 & 2.6640 & 3.3137 \\ \hline
			20 & 2.3412 & 3.1837 \\ \hline
			23 & 1.8415 & 2.9577 \\ \hline
			26 & 1.3656 & 2.7210 \\ \hline
			30 & 0.8264 & 2.4280 \\ \hline
			34 & 0.4124 & 2.1736 \\ \hline
			38 & 0.1231 & 1.9527 \\ \hline
		\end{tabular}
		\caption*{NewHope, Hamming}
	\end{minipage}
	\begin{minipage}{0.48\textwidth}
		\centering
		\begin{tabular}{|c|c|c|}
			\hline Q & $R_{GV}$ & $R_H$ \\ \hline
			2 & 1.0000 & 1.0000 \\ \hline
			3 & 1.5850 & 1.5850 \\ \hline
			4 & 2.0000 & 2.0000 \\ \hline
			5 & 2.3219 & 2.3219 \\ \hline
			6 & 2.5645 & 2.5742 \\ \hline
			7 & 2.7868 & 2.7966 \\ \hline
			8 & 2.9448 & 2.9703 \\ \hline
			9 & 3.0989 & 3.1314 \\ \hline
			10 & 3.2070 & 3.2587 \\ \hline
			12 & 3.3331 & 3.4429 \\ \hline
			14 & 3.3518 & 3.5442 \\ \hline
			16 & 3.3163 & 3.5966 \\ \hline
			18 & 3.2367 & 3.6080 \\ \hline
			20 & 3.1565 & 3.6074 \\ \hline
			23 & 3.0474 & 3.5969 \\ \hline
			26 & 2.9621 & 3.5871 \\ \hline
			30 & 2.8897 & 3.5867 \\ \hline
			34 & 2.8525 & 3.5995 \\ \hline
			38 & 2.8344 & 3.6182 \\ \hline
		\end{tabular}
		\caption*{NewHope, Li}
	\end{minipage}
\end{table}

Видно, что при маленьких $Q$ и $t$ разница между метриками Хэмминга и Ли небольшая или вообще отсутствует, при увеличении же этих параметров Хэмминг начинает резко убывать, а Ли почти не уменьшается (за счёт гораздо меньшего объёма шаров, возникающих вокруг кодовых слов).

Также занимательно, что вероятность возникновения отличия на $\geq\!2$ в метрике Ли перестаёт быть пренебрежительно мала лишь к моменту, когда вероятность отличия на $1$ становится огромна (порядка $\frac14$ и выше).

Таким образом, с точки зрения кол-ва исправляемых ошибок при малых и средних значениях $Q$ различий между метриками Хэмминга и Ли вообще не наблюдается (параметр $t$ при данном DFR совпадает).
